% ---- page & fonts -----------------------------------------------------------
\usepackage[T1]{fontenc}
\usepackage{lmodern}
\usepackage[margin=2.5cm]{geometry}
\usepackage{enumitem}
\usepackage[dvipsnames]{xcolor}

% ---- math -------------------------------------------------------------------
\usepackage{amsmath,amssymb,amsthm,mathtools,bm}

% ---- links (cleveref must load after hyperref and before \newtheorem) ----------
\usepackage[colorlinks=true,linkcolor=NavyBlue,citecolor=NavyBlue,urlcolor=NavyBlue]{hyperref}
\usepackage[capitalise,noabbrev]{cleveref}

% ---- theorem environments (one shared counter per chapter) -------------------
\theoremstyle{plain}
\newtheorem{theorem}{Theorem}[chapter]
\newtheorem{lemma}[theorem]{Lemma}
\newtheorem{proposition}[theorem]{Proposition}
\newtheorem{corollary}[theorem]{Corollary}
\theoremstyle{definition}
\newtheorem{definition}[theorem]{Definition}
\newtheorem{example}[theorem]{Example}
\theoremstyle{remark}
\newtheorem{remark}[theorem]{Remark}

% ---- entry metadata ---------------------------------------------------------
% \source{...}: where the result comes from (lecture, book section, paper).
\newcommand{\source}[1]{{\footnotesize\upshape\textcolor{gray}{Source: #1}}\par\smallskip}
% \derived{YYYY-MM-DD}: date I reproduced the full proof closed-book (mastery-gate standard).
\newcommand{\derived}[1]{\par\hfill{\footnotesize\textsf{$\checkmark$ derived closed-book: #1}}\par}
% \pending: proof not yet written / not yet reproducible without notes.
\newcommand{\pending}{\par\hfill{\footnotesize\textsf{\textcolor{BrickRed}{proof pending}}}\par}

% ---- notation macros (documented in chapters/00-notation.tex) ---------------
\newcommand{\R}{\mathbb{R}}
\newcommand{\C}{\mathbb{C}}
\newcommand{\N}{\mathbb{N}}
\newcommand{\tp}{^{\top}}
\DeclarePairedDelimiter{\norm}{\lVert}{\rVert}
\DeclarePairedDelimiter{\abs}{\lvert}{\rvert}
\DeclarePairedDelimiterX{\inner}[2]{\langle}{\rangle}{#1, #2}
\DeclareMathOperator{\rank}{rank}
\DeclareMathOperator{\tr}{tr}
\DeclareMathOperator{\diag}{diag}
\DeclareMathOperator{\colsp}{C}
\DeclareMathOperator{\nullsp}{N}
\DeclareMathOperator*{\argmin}{arg\,min}
\DeclareMathOperator*{\argmax}{arg\,max}
